\section{Solutions for Chapter~\ref{ch:haagerup}}

\paragraph{Exercise~\ref{ex:H3-fusion}.}
\phantomsection
\label{sol:ex:H3-fusion}
The invertible factors give
$A_gA_h=\alpha^{g+h}=A_{g+h}$ and
$A_gR_h=\alpha^{g+h}\rho=R_{g+h}$.
The relation $\rho\alpha^h=\alpha^{-h}\rho$ gives
$R_gA_h=\alpha^{g-h}\rho$.  It also gives
\[
 R_gR_h=\alpha^{g-h}\rho^2
 =A_{g-h}+\sum_kR_{g-h+k}
 =A_{g-h}+\sum_kR_k.
\]
The product of $A_g$ with an invertible object contains the unit
exactly for $A_gA_{-g}$, and $R_gR_h$ contains the unit
exactly when $g=h$, proving the duality rule.  Finally,
$\alpha\rho=R_1$ and $\rho\alpha=R_2$ are distinct simples.
A braiding would make them isomorphic, so no braiding exists on
this six-object category.



\paragraph{Exercise~\ref{ex:H3-F}.}
\phantomsection
\label{sol:ex:H3-F}
Put $r=\sqrt{13}$. From the radicals,
\[
 y_-+y_+=\frac{5-r}{6},\qquad
 y_-y_+=\frac{(5-r)^2-6(1+r)}{144}
          =\frac{2-r}{9}=\frac{x}{3}.
\]
Consequently
\[
 xy_-+y_-y_++y_+x
 =\frac{(2-r)(7-r)}{18}
 =\frac{3-r}{2}=-d^{-1}.
\]
Together with $x+y_-+y_+=-d^{-1}$, this gives
$x^2+y_-^2+y_+^2=d^{-2}+2d^{-1}=1-d^{-1}$.
The first row of the full matrix has norm squared
$d^{-2}+3d^{-1}=1$. Its inner product with any lower row is
$d^{-1/2}(d^{-1}+x+y_-+y_+)=0$. Each lower row has norm
squared $d^{-1}+x^2+y_-^2+y_+^2=1$, and two distinct lower
rows have inner product $d^{-1}+xy_-+y_-y_++y_+x=0$.
Thus all rows are orthonormal. This checks the inverse of one
recoupling map. The pentagon also compares compositions of
different blocks, for every admissible choice of external labels;
those equations have not been checked by this calculation.



\paragraph{Exercise~\ref{ex:H3-P}.}
\phantomsection
\label{sol:ex:H3-P}
Let $S_0=\sum_gY_{A_g}$ and $S_1=\sum_gY_{R_g}$.  Summing
the fusion rules gives
\[
 S_0^2=3S_0,\quad S_0S_1=S_1S_0=3S_1,\quad
 S_1^2=3S_0+9S_1.
\]
For $N=S_0+dS_1$, it follows that
\[
 N^2=3(1+d^2)S_0+(6d+9d^2)S_1=3sN,
\]
where $6d+9d^2=3sd$ follows from $d^2=1+3d$.
Therefore $P_+=N/(3s)$ is idempotent.  Also
\[
 Y_\rho S_0=S_1,\qquad
 Y_\rho S_1=S_0+3S_1,
\]
so $Y_\rho N=dS_0+(1+3d)S_1=dN$.
No eigenvectors, associator entries, or numerical approximations
are needed for this ring identity.



\paragraph{Exercise~\ref{ex:H3-characteristic}.}
\phantomsection
\label{sol:ex:H3-characteristic}
The block-constant subspace gives characteristic polynomial
$x^2-3x-1$.  On the two copies of the sum-zero subspace, the
matrix exchanges the blocks, giving $+1$ and $-1$ each twice.
Thus
\[
 \det(xI-N_\rho)=(x^2-3x-1)(x^2-1)^2.
\]
Taking traces of powers gives \eqref{eq:H3-path-count}.
The recursion yields
$C_2=11,C_3=36,C_4=119,C_5=393,C_6=1298$; adding the
even-length contribution four gives $1302$ paths at $L=6$.



\paragraph{Exercise~\ref{ex:H3-neutral}.}
\phantomsection
\label{sol:ex:H3-neutral}
On $\C^3\oplus\C^2\oplus(\C^2)^{\oplus4}$ with its standard
positive form, take
\begin{align*}
 Y_\alpha&=I_3\oplus I_2\oplus
     \begin{pmatrix}\omega&0\\0&\omega^2\end{pmatrix}^{\oplus4},\\
 Y_\rho&=dI_3\oplus(-d^{-1}I_2)\oplus
     \begin{pmatrix}0&1\\1&0\end{pmatrix}^{\oplus4}.
\end{align*}
Each block obeys the fusion rules, and
$Y_\alpha^\dagger=Y_{\alpha^2}$, $Y_\rho^\dagger=Y_\rho$.
This is a unitary $*$-representation of dimension
$3+2+4\cdot2=13$; it does not require the noninvertible operator
$Y_\rho$ itself to be a unitary matrix.
The neutral projector $E$ has rank five, since it is the identity
on both kinds of scalar representation and zero on every charged
doublet.  The ranks of $P_+$, $P_-$, and $P_{\rm ch}$ are
$3,2,8$, respectively.  There are five $\alpha$-neutral states
but only three fully Haagerup-invariant states.  Neutrality under
the invertible subgroup is therefore a strictly weaker condition.


